% =============================================================================
%  vsim_code_request.tex
%  Implementation Request — Bridge Layer
%  CR-2026-01 · VSIM Internal · v5.2.0-internal
% =============================================================================
\documentclass[11pt, a4paper]{article}
\usepackage{vsim_common}

\hypersetup{
  pdftitle={Code Request CR-2026-01 --- Bridge Layer Implementation},
  pdfsubject={VSIM Internal --- v5.2.0-internal}
}

% =============================================================================
\begin{document}
% =============================================================================

\thispagestyle{empty}

\begin{center}
  \eyebrow{VSIM Internal \textperiodcentered\ Code Request}\\[8pt]
  {\LARGE\sffamily\bfseries Implementation Request --- Bridge Layer}\\[6pt]
  {\normalsize\sffamily
    Identity Entropy Output Layer \quad\textperiodcentered\quad
    Formation Context Wiring \quad\textperiodcentered\quad
    Coalescence Tracking}
\end{center}

\vspace{10pt}
\noindent
\begin{tabular}{@{}L{3cm} L{4.8cm} L{3cm} L{4.2cm}@{}}
  \code{Doc ID}    & \texttt{CR-2026-01}                    &
  \code{Status}    & Open --- pending implementation \\
  \code{Target}    & \texttt{v5.2.0-internal}               &
  \code{Audience}  & VSEPR-SIM / VSIM developers \\
  \code{Raised}    & 2026-05-31                             &
  \code{Closed}    & --- \\
\end{tabular}

\vspace{8pt}
\begin{infobox}
  \textbf{Ref:} Theoretical--Simulation Bridge Project
  (\texttt{vsim\_bridge\_doc.tex}) --- read that first.
  This document assumes familiarity with \S\,1--10 of the bridge doc
  and does not re-derive theory.
\end{infobox}

\noindent\rule{\linewidth}{2pt}
\vspace{6pt}

\tableofcontents
\newpage

% =============================================================================
\section{Summary}
% =============================================================================

The bridge doc defines what the identity entropy layer should do. This document converts
that into ten concrete implementation tasks with priorities, interface contracts, open
decisions, and acceptance criteria. Tasks are grouped into three priority tiers. \pzero\
tasks are blocking for v5.2.0. \pone\ tasks complete the feature. \ptwo\ tasks can defer
to a follow-on milestone.

\vspace{8pt}
\begin{center}
\small\sffamily
\begin{tabular}{@{} C{2.5cm} C{2.5cm} C{2.5cm} C{2.5cm} C{2.5cm} @{}}
  \toprule
  \textbf{Total tasks} & \textbf{P0 --- blocking} & \textbf{P1 --- required} &
  \textbf{P2 --- deferrable} & \textbf{Open decisions} \\
  \midrule
  10 & 3 & 5 & 2 & 5 \\
  \bottomrule
\end{tabular}
\end{center}

% =============================================================================
\section{Implementation Tasks}
% =============================================================================

\begin{warnbox}
  \textbf{Dependencies:} CR-03 and CR-04 must land before CR-05 through CR-09 can
  be meaningfully tested. CR-01 must land before CR-07. CR-06 has a hard dependency
  on DEC-05 being resolved.
\end{warnbox}

\vspace{6pt}
\small
\begin{longtable}{@{} L{1.1cm} L{7.8cm} L{1.2cm} C{1.1cm} C{1.4cm} @{}}
  \rowcolor{skybg}
  \textbf{ID} & \textbf{Task} & \textbf{Ref} & \textbf{Pri.} & \textbf{Status} \\
  \toprule
  \endhead
  \bottomrule
  \endfoot

  \texttt{CR-01} &
    \textbf{Parse \code{[formation]} and \code{[formation.initial\_state]} into
    \code{FormationContext}.}
    Struct skeleton in \code{vsim/identity\_entropy.hpp}.
    Fields: \code{method}, \code{history}, \code{formation\_time},
    \code{temperature\_K}, \code{pressure\_Pa}, \code{phase},
    \code{defect\_state}, \code{grain\_state}, \code{density}.
    Absent sub-fields default to \code{"unknown"} / zero, not error.
    & \S\,4, \S\,10 & \pzero & \spend \\[4pt]

  \texttt{CR-02} &
    \textbf{Formation validation pass --- mark runs with absent or incomplete
    formation as \code{incomplete} in the validation layer.}
    Applies to simulations with material types: solid, surface, gas, pipe, powder,
    structural. Missing \code{[formation]} block produces a validation warning,
    not a hard failure. Runs must still execute.
    & \S\,4 & \pzero & \spend \\[4pt]

  \texttt{CR-03} &
    \textbf{Parse \code{[outputs.identity\_entropy]} config block and wire into
    the export pipeline.}
    Fields: \code{enabled}, \code{path}, \code{format}, \code{precision},
    \code{attach\_to\_viewer}, \code{write\_per\_frame}, \code{write\_summary},
    \code{variables}.
    If \code{enabled = false} the block is parsed but no writer is instantiated.
    Schema decision DEC-01 must be resolved before this lands.
    & \S\,5 & \pzero & \spend \\[4pt]

  \texttt{CR-04} &
    \textbf{Implement \code{IEWriter} --- header, per-frame write, summary, close.}
    Full implementation stub in \code{vsim/identity\_entropy.cpp} (bridge doc \S\,10).
    \code{write(IERecord\&)} must call \code{sync()} before writing.
    Writer must not throw after \code{close()} is called.
    Buffer strategy to be confirmed per DEC-04.
    & \S\,5, \S\,10 & \pzero & \spend \\[4pt]

  \texttt{CR-05} &
    \textbf{Populate all 12 \code{IERecord} fields per simulation frame.}
    \code{J}, \code{dataloss}, \code{recoverability} have defined derivation rules.
    \code{dataentropy}, \code{hidden\_activity}, and \code{projection\_loss}
    require physics model input --- see DEC-02.
    \code{state\_quota} maps to $Q_n$. \code{entropy\_flux} is the frame-delta of
    identity entropy. Fields not yet computable emit \code{0.0} with flag bit set.
    & \S\,6 & \pone & \spend \\[4pt]

  \texttt{CR-06} &
    \textbf{Wire \code{coalescence\_gate()} into the bonding/clustering event system.}
    When the gate fires, emit nonzero \code{coalescence\_loss} and \code{residual}
    into the \code{IERecord} for that frame. Gate is \emph{read-only} --- it must
    not mutate particle state on fire. \code{coalescence\_residual()} provides the
    scalar loss value. See DEC-05.
    & \S\,7, \S\,10 & \pone & \spend \\[4pt]

  \texttt{CR-07} &
    \textbf{Implement \code{FormationContext::formation\_memory()} and wire its
    result into \code{IERecord.formation\_memory} each frame.}
    The heuristic in bridge doc \S\,10 scores six declared descriptors.
    A real implementation should degrade $M_f$ over time as the system
    thermalizes. The heuristic is the starting floor. Thermalization model is
    out of scope for P1; static heuristic is sufficient to unblock.
    & \S\,8, \S\,10 & \pone & \spend \\[4pt]

  \texttt{CR-08} &
    \textbf{Add \code{"identity\_entropy"} as a valid entry in
    \code{[observe].track}.}
    When present, the runtime activates identity entropy tracking.
    If the entry is present but \code{[outputs.identity\_entropy]} is absent or
    \code{enabled = false}, emit a validation warning (not a failure).
    & \S\,9 & \pone & \spend \\[4pt]

  \texttt{CR-09} &
    \textbf{Per-frame vs summary-only write modes, controlled by
    \code{write\_per\_frame} and \code{write\_summary} flags.}
    \code{write\_per\_frame = false} + \code{write\_summary = true} produces
    a summary-only \code{.ie} file.
    Both false is a no-op that still creates the file with just the header.
    Both true is the default full output.
    & \S\,5 & \pone & \spend \\[4pt]

  \texttt{CR-10} &
    \textbf{Viewer overlay wiring when \code{attach\_to\_viewer = true}.}
    The \code{.ie} file path should be registered with the viewer's overlay system.
    Exact overlay API is viewer-team's call --- this task is just the registration
    handoff from the export pipeline. Can ship without this; viewer overlay is not
    required for the output layer to be useful.
    & \S\,5 & \ptwo & \spend \\

\end{longtable}
\normalsize

% =============================================================================
\section{Interface Contracts}
% =============================================================================

These constraints must not be violated regardless of how tasks are subdivided
internally. They follow directly from the bridge doc's core principle.

\vspace{6pt}
\begin{infobox}
\small
\begin{description}[leftmargin=2.4cm, style=sameline, font=\ttfamily\small]
  \item[IC-01] Identity entropy variables (\code{J}, \code{dataloss}, all
    \code{IERecord} fields) must never be written back into mutable simulation
    state. They are derived output only.
  \item[IC-02] The \code{coalescence\_gate()} function is read-only. A gate fire
    produces a sidecar record. It does not create, remove, or modify particles.
  \item[IC-03] Render state must not mutate truth state. The viewer overlay (CR-10)
    reads from the \code{.ie} file; it does not feed back into the simulation kernel.
  \item[IC-04] \code{FormationContext} is a runtime-only object. It must not be
    serialised into \code{.dynx} or any ground-truth particle record.
  \item[IC-05] \code{IEWriter::write(IERecord\&)} must call \code{sync()} before
    writing. Callers must not pre-call \code{sync()} and assume the writer will skip
    it --- double-sync is safe; missing sync is not.
  \item[IC-06] A missing or disabled \code{[outputs.identity\_entropy]} block must
    not degrade simulation performance. The code path must be a no-op when disabled.
  \item[IC-07] \code{J} must be in the range $[0.0, 1.0]$ at all times. Clamp on
    computation, not on write. Values outside range indicate a model error.
\end{description}
\end{infobox}

% =============================================================================
\section{Open Decisions}
% =============================================================================

These must be resolved before the tasks that depend on them can land.

\vspace{6pt}
\begin{description}[leftmargin=2cm, style=nextline, font=\sffamily\bfseries\small]

  \item[DEC-01 \hfill \stbd]
    Schema bracket: \code{[outputs.identity\_entropy]} or
    \code{[export.identity\_entropy]}?\\
    {\small The bridge doc deliberately deferred this. Pick one and stick to it.
    Blocks CR-03. If the rest of the export system uses \code{[export.*]}
    sub-tables, consistency suggests folding in. If \code{[outputs.*]} is the
    new direction, keep it separate.}

  \item[DEC-02 \hfill \stbd]
    How are \code{dataentropy}, \code{hidden\_activity}, and
    \code{projection\_loss} computed?\\
    {\small These require a model of the projection operator from identity-space to
    observable state. The bridge doc defines the concept but not the formula. Until
    resolved, CR-05 emits \code{0.0} with flag bit \texttt{0x1} set. Blocks CR-05
    (full), CR-06 (partial).}

  \item[DEC-03 \hfill \stbd]
    \code{lineage\_depth} tracking strategy: per-object chain or global lineage
    registry?\\
    {\small Each object could carry its own depth counter (simple, cheap, no global
    state), or the runtime could maintain a lineage map keyed on identity ID
    (accurate for re-coalesced objects, more memory). Blocks CR-05 for
    \code{lineage\_depth} population.}

  \item[DEC-04 \hfill \stbd]
    \code{IEWriter} buffer strategy: stream per frame or accumulate $N$ frames
    then flush?\\
    {\small Streaming per frame is simpler and crash-safe. Buffering $N$ frames is
    faster for large simulations. The summary pass already buffers all records in
    memory --- if memory is a concern the summary could be a running accumulator.
    Blocks CR-04 finalisation.}

  \item[DEC-05 \hfill \stbd]
    Definition of the projection operator $A_n$ for coalescence residual
    computation.\\
    {\small \code{coalescence\_residual()} takes caller-supplied
    \code{projected\_identity} vector. Who computes the projection, and how?
    The bridge doc defines
    $R_{\text{coal}}^{(n)} = C_n - A_n(I^{\star}_{(n+1)})$
    but $A_n$ is not yet specified. Blocks CR-06.}

\end{description}

% =============================================================================
\section{Acceptance Criteria}
% =============================================================================

A milestone is shippable when all \pzero\ and \pone\ criteria below are checked off.
\ptwo\ criteria are tracked but not required for v5.2.0.

\vspace{6pt}
\small
\begin{longtable}{@{} L{1.3cm} C{0.7cm} L{10.5cm} C{1.1cm} @{}}
  \rowcolor{skybg}
  \textbf{ID} & \textbf{Tick} & \textbf{Criterion} & \textbf{Priority} \\
  \toprule
  \endhead
  \bottomrule
  \endfoot

  AC-01 & \cb &
    A surface deposition run with \code{[outputs.identity\_entropy] enabled = true}
    produces a valid \code{.ie} file at the declared path.
    & \pzero \\

  AC-02 & \cb &
    The \code{.ie} file contains a header row, at least one data row per frame,
    and a \code{\# SUMMARY} block with \code{rows}, \code{mean\_J}, \code{min\_J},
    and \code{coal\_events}.
    & \pzero \\

  AC-03 & \cb &
    \code{J} is within $[0.0, 1.0]$ for all rows. \code{dataloss} equals
    $1 - J$ to within \texttt{float64} tolerance.
    \code{recoverability} equals $1 - \text{projection\_loss}$.
    & \pone \\

  AC-04 & \cb &
    A run with all six formation descriptors declared produces
    \code{formation\_memory = 1.0} on frame 0.
    A run with no formation block produces \code{formation\_memory = 0.0}.
    & \pone \\

  AC-05 & \cb &
    A run without a \code{[formation]} block executes successfully and
    the validation report contains a \code{formation: incomplete} warning.
    No hard failure is emitted.
    & \pzero \\

  AC-06 & \cb &
    When a coalescence event fires, the corresponding frame row in the \code{.ie}
    file has \code{coalescence\_loss > 0.0}. Frames without coalescence events
    have \code{coalescence\_loss = 0.0}.
    & \pone \\

  AC-07 & \cb &
    \code{[outputs.identity\_entropy] enabled = false} produces no \code{.ie}
    file and no measurable performance overhead versus a run without the block.
    & \pzero \\

  AC-08 & \cb &
    A script with \code{"identity\_entropy"} in \code{[observe].track} but no
    \code{[outputs.identity\_entropy]} block emits a validation warning and
    runs without error.
    & \pone \\

  AC-09 & \cb &
    \code{write\_per\_frame = false} + \code{write\_summary = true} produces a
    \code{.ie} file with only the header and \code{\# SUMMARY} block ---
    no data rows.
    & \pone \\

  AC-10 & \cb &
    With \code{attach\_to\_viewer = true}, the \code{.ie} file path is registered
    with the viewer overlay system. The viewer can display per-identity entropy
    values without the simulation re-running.
    & \ptwo \\

\end{longtable}
\normalsize

% =============================================================================
\section{File Inventory}
% =============================================================================

New files introduced by this request. Existing files should not require structural
changes beyond adding new parser branches and wiring calls.

\vspace{6pt}
\small
\begin{tabular}{@{} L{4.5cm} L{2.5cm} L{5.0cm} L{2.5cm} @{}}
  \rowcolor{skybg}
  \textbf{File} & \textbf{Type} & \textbf{Contents} & \textbf{Tasks} \\
  \toprule

  \code{vsim/identity\_entropy.hpp} &
    New header &
    \code{FormationContext}, \code{IERecord}, \code{CoalescenceInput},
    \code{coalescence\_gate()}, \code{coalescence\_residual()}, \code{IEWriter}.
    Skeleton in bridge doc \S\,10. &
    CR-01, 04, 06, 07 \\[4pt]

  \code{vsim/identity\_entropy.cpp} &
    New TU &
    Implementations of all non-inline functions from the header.
    Skeleton in bridge doc \S\,10. &
    CR-04, 06, 07 \\[4pt]

  \code{vsim/parser/} \code{formation\_parser.*} &
    New parser module &
    Reads \code{[formation]} and \code{[formation.initial\_state]} into
    \code{FormationContext}. Follows existing parser conventions. &
    CR-01 \\[4pt]

  \code{vsim/export/} \code{ie\_export.*} &
    New export module &
    Reads \code{[outputs.identity\_entropy]} config, instantiates \code{IEWriter},
    wires frame callbacks into existing export pipeline alongside \code{.dynx},
    analysis JSON, etc. &
    CR-03, 05, 08, 09 \\[4pt]

  \code{vsim/validation/} \code{formation\_check.*} &
    New validation check &
    Emits \code{formation: incomplete} warning for runs missing
    \code{[formation]}. Runs after config parse, before runtime init. &
    CR-02 \\

  \bottomrule
\end{tabular}
\normalsize

\vspace{12pt}
\begin{warnbox}
  \textbf{Path convention:}
  Module paths above follow the existing VSIM source tree convention and are
  illustrative. Actual paths should follow whatever the team has established for v5.x.
  Do not bikeshed the filenames; do bikeshed the interfaces.
\end{warnbox}

% =============================================================================
\vspace{16pt}
\noindent\rule{\linewidth}{0.5pt}
{\small\sffamily\textcolor{gray}{CR-2026-01 \textperiodcentered\
  VSIM Internal \textperiodcentered\ v5.2.0-internal
  \textperiodcentered\ Ref: vsim\_bridge\_doc.tex}}

\end{document}
